\section{Related Work}

\paragraph{Learning from model weights.}
Weight-space learning covers work on weight geometry, on learned representations of parameters, and on generating new parameters \citep{han2026survey}. Our question is a narrower one: predicting measured performance changes from the weight difference between successive checkpoints. W2T is a close precedent for the representation side of this problem. It canonicalizes LoRA updates through QR decomposition and SVD, tokenizes the resulting components, and processes them with rank-level and position-level Transformers. On its ARC-LoRA benchmark, W2T reports test MAE 0.32 and Spearman 95.35\% \citep{w2t2026}. Those numbers concern LoRA performance prediction on that benchmark. We cannot read off their performance on merged edits, multi-edit trajectories, or the benchmarks in our proposed study. We therefore treat canonicalized low-rank features as a candidate representation, not as evidence that the same accuracy will carry over.

Several neighboring lines predict performance by running the model first. Delta Activations represents fine-tuned models by measuring how their internal activations change relative to a base model \citep{deltaactivations2025}. Fast probing reads intermediate hidden activations to predict downstream performance. It reports AUROC above 0.75 and cuts evaluation latency from about an hour to three minutes on OLMo3-7B-base \citep{fastprobing2026}. TuneAhead pairs dataset descriptors with short-probe features and reports a 1.47 percentage-point RMSE over 370 held-out Qwen2.5-7B-Instruct runs \citep{tuneahead2026}. These approaches make a cost-aware comparison worth doing, but their inputs include activation or probe information. Our main setting uses weight deltas and their summaries, with no extra forward passes. Any accuracy comparison across these lines therefore has to account for different inputs and compute budgets.

\paragraph{Edit trajectories and predictive representations.}
Weight changes may carry information about an edit's effect. The effect of a later edit, though, need not be independent of the edits before it. Work on task-vector composition finds that interference depends on training-format serialization: a trained wrapper can mask interference that appears under an untrained chat template \citep{taskvectors2026interfere}. This motivates recording evaluation formats and modeling trajectory history, and it cautions against treating a merge as context-free addition. A separate study across CIFAR training and Pythia checkpoints reports localized, time-varying pockets of predictable short-horizon dynamics in bulk parameters \citep{predictability2026dynamics}. Leap+Verify proposes speculative future-weight predictions and accepts them only after validation against a held-out loss criterion \citep{leapverify2026}. That is a proposal about prediction with verification. It is not evidence that our performance predictor will reduce training time.

We test sequence models because each observed edit arrives in context. An edit can change the starting point for the next one. Mamba-3 pairs a more expressive state-space recurrence with complex-valued state updates and a multi-input, multi-output formulation \citep{mamba3}. That description is architectural, and it does not show an advantage on model-edit trajectories. Our H2 comparison asks whether a stateful sequence predictor improves held-out performance estimates over a stateless per-edit model. A separate JEPA-style model will learn to predict a representation of the next weight state from the prior state and an edit descriptor. The JEPAWG work on physical parameterized systems omits the VICReg covariance term from its main-text configuration, which allows a lower-dimensional latent representation \citep{jepawg2026physics}. For predictor initialization, ResidualPred is reported to bias transformer attention toward identity. In the reported tabular and image experiments it also improves effective rank and downstream accuracy \citep{drivedecay2026}. Those findings inform our design choices, but neither one shows transfer across language-model edit families. JEPA and hidden-Markov-model correspondences hold at different levels: exact sequence-level equivalence needs additional stated conditions, and standard JEPA training alone does not supply them \citep{jepahmm2026}.

\paragraph{Alignment edits and evaluation limits.}
Refusal-related weight edits make a useful stress test, because refusal can change while task capability changes too. Heretic combines configurable modifiers with Optuna TPE optimization, and its default modifier is Arbitrary-Rank Ablation \citep{heretic2026}. We plan to compare this family with classic directional ablation. Capability goes into the measurements; we do not assume it survives. A separate analysis of Heretic's default refusal score notes that it matches any of 33 case-insensitive strings within the first 100 response tokens, and the score alone does not measure harmfulness \citep{pksharma2026heretic}. This motivates reporting refusal measures next to capability and distributional change, and reading keyword-based rates narrowly.

Prior work also illustrates why neither safety nor collateral effects should be inferred from one intervention or proxy. Surgical Refusal Ablation orthogonalizes a refusal vector against protected concept atoms. Across five models it reports low refusal rates alongside measured perplexity and first-token KL changes, and those measurements do not establish universal preservation of task accuracy \citep{surgicalrefusal2026}. At larger scale, directional ablation on GLM-5.3-Flash (320B) left measurable refusal after every tested edit, including tested subspaces of ranks 1--12 \citep{singleattack2026moe}. That is evidence about the edits that were tested. It does not show that directional ablation cannot change refusal in general. Our alignment sub-study treats outcomes like these as reasons to measure several endpoints, and we do not take them as predicted results for Qwen. Across both the weight-space and alignment literatures, then, our contribution is an empirical test of whether changes can be predicted from deltas and trajectories. We keep the boundaries of the source claims in view and do not assume that prediction can substitute for evaluation.
